\documentclass{amsart}
% For dealing with references we use the comment environment
\usepackage{verbatim}
\newenvironment{reference}{\comment}{\endcomment}
%\newenvironment{reference}{}{}
\newenvironment{slogan}{\comment}{\endcomment}
\newenvironment{history}{\comment}{\endcomment}

% For commutative diagrams we use Xy-pic
\usepackage[all]{xy}

% We use 2cell for 2-commutative diagrams.
\xyoption{2cell}
\UseAllTwocells

% We use multicol for the list of chapters between chapters
\usepackage{multicol}

% This is generally recommended for better output
\usepackage{lmodern}
\usepackage[T1]{fontenc}

% For cross-file-references

% Package for hypertext links:
\usepackage{hyperref}

% For any local file, say "hello.tex" you want to link to please

% Theorem environments.
%
\theoremstyle{plain}
\newtheorem{theorem}[subsection]{Theorem}
\newtheorem{proposition}[subsection]{Proposition}
\newtheorem{lemma}[subsection]{Lemma}

\theoremstyle{definition}
\newtheorem{definition}[subsection]{Definition}
\newtheorem{example}[subsection]{Example}
\newtheorem{exercise}[subsection]{Exercise}
\newtheorem{situation}[subsection]{Situation}

\theoremstyle{remark}
\newtheorem{remark}[subsection]{Remark}
\newtheorem{remarks}[subsection]{Remarks}

\numberwithin{equation}{subsection}

% Macros
%
\def\lim{\mathop{\mathrm{lim}}\nolimits}
\def\colim{\mathop{\mathrm{colim}}\nolimits}
\def\Spec{\mathop{\mathrm{Spec}}}
\def\Hom{\mathop{\mathrm{Hom}}\nolimits}
\def\Ext{\mathop{\mathrm{Ext}}\nolimits}
\def\SheafHom{\mathop{\mathcal{H}\!\mathit{om}}\nolimits}
\def\SheafExt{\mathop{\mathcal{E}\!\mathit{xt}}\nolimits}
\def\Sch{\mathit{Sch}}
\def\Mor{\mathop{\mathrm{Mor}}\nolimits}
\def\Ob{\mathop{\mathrm{Ob}}\nolimits}
\def\Sh{\mathop{\mathit{Sh}}\nolimits}
\def\NL{\mathop{N\!L}\nolimits}
\def\CH{\mathop{\mathrm{CH}}\nolimits}
\def\proetale{{pro\text{-}\acute{e}tale}}
\def\etale{{\acute{e}tale}}
\def\QCoh{\mathit{QCoh}}
\def\Ker{\mathop{\mathrm{Ker}}}
\def\Im{\mathop{\mathrm{Im}}}
\def\Coker{\mathop{\mathrm{Coker}}}
\def\Coim{\mathop{\mathrm{Coim}}}

% Boxtimes
%
\DeclareMathSymbol{\boxtimes}{\mathbin}{AMSa}{"02}

%
% Macros for moduli stacks/spaces
%
\def\QCohstack{\mathcal{QC}\!\mathit{oh}}
\def\Cohstack{\mathcal{C}\!\mathit{oh}}
\def\Spacesstack{\mathcal{S}\!\mathit{paces}}
\def\Quotfunctor{\mathrm{Quot}}
\def\Hilbfunctor{\mathrm{Hilb}}
\def\Curvesstack{\mathcal{C}\!\mathit{urves}}
\def\Polarizedstack{\mathcal{P}\!\mathit{olarized}}
\def\Complexesstack{\mathcal{C}\!\mathit{omplexes}}
% \Pic is the operator that assigns to X its picard group, usage \Pic(X)
% \Picardstack_{X/B} denotes the Picard stack of X over B
% \Picardfunctor_{X/B} denotes the Picard functor of X over B
\def\Pic{\mathop{\mathrm{Pic}}\nolimits}
\def\Picardstack{\mathcal{P}\!\mathit{ic}}
\def\Picardfunctor{\mathrm{Pic}}
\def\Deformationcategory{\mathcal{D}\!\mathit{ef}}

\setlength{\emergencystretch}{3em}
\begin{document}
\title[Stacks Project, Tag 0AA8]{254 --- Tor-independent base change for the right adjoint of proper pushforward}
\maketitle
\noindent Copyright (C) 2005--2025 Johan de Jong. Stacks Project authors. Source: \href{https://stacks.math.columbia.edu/tag/0AA8}{Tag 0AA8}.

\noindent Editorial difficulty note (not author text). Difficulty H2: The argument tests the canonical base-change map against a perfect generator and compares two global derived Hom complexes. It separately controls perfect and pseudo-coherent pushforward hypotheses, then proves the constructed identification is exactly the specified natural transformation through a large adjunction diagram. Identifying the map, rather than only isomorphic objects, is a core obligation..

\noindent Editorial provenance note (not author text). History: excerpt from revision \texttt{a0444\allowbreak 6e57e\allowbreak c1fbc\allowbreak 252a8\allowbreak 71afc\allowbreak ec775\allowbreak 2fb28\allowbreak 07b14}, duality.tex, lines 1418--1606. Reference presentation was adapted; original-author context and core support blocks were retained or added. The mathematical wording is unchanged. Licensed under GNU FDL 1.2 or later; no invariant sections or cover texts. See ../licenses/COPYING. Context blocks reproduce author definitions and supporting results; they are not additional counted problems.

\medskip\noindent
We often want to know whether the right adjoints to pushforward commutes
with base change. Thus we consider a cartesian square
\begin{equation}
\label{equation-base-change}
\vcenter{
\xymatrix{
X' \ar[r]_{g'} \ar[d]_{f'} & X \ar[d]^f \\
Y' \ar[r]^g & Y
}
}
\end{equation}
of quasi-compact and quasi-separated schemes.
Denote
$$
a  : D_\QCoh(\mathcal{O}_Y) \to D_\QCoh(\mathcal{O}_X)
\quad\text{and}\quad
a'  : D_\QCoh(\mathcal{O}_{Y'}) \to D_\QCoh(\mathcal{O}_{X'})
$$
the right adjoints to $Rf_*$ and $Rf'_*$
(Lemma \href{https://stacks.math.columbia.edu/tag/0A9E}{lemma twisted inverse image (Tag 0A9E)}).
Consider the base change map of
Cohomology, Remark \href{https://stacks.math.columbia.edu/tag/08HY}{remark base change (Tag 08HY)}.
It induces a transformation of functors
$$
Lg^* \circ Rf_* \longrightarrow Rf'_* \circ L(g')^*
$$
on derived categories of sheaves with quasi-coherent cohomology.
Hence a transformation between the right adjoints in the opposite direction
$$
a \circ Rg_* \longleftarrow Rg'_* \circ a'
$$

\noindent
The map (\href{https://stacks.math.columbia.edu/tag/0A9L}{equation sheafy (Tag 0A9L)}) is a special case of a base change map.
Namely, suppose that we have a cartesian diagram
$$
\xymatrix{
X' \ar[r]_{g'} \ar[d]_{f'} & X \ar[d]^f \\
Y' \ar[r]^g & Y
}
$$
of quasi-compact and quasi-separated schemes, i.e., a diagram as in
(\ref{equation-base-change}). Assume $f$ and $g$ are {\bf Tor independent}.
Then we can consider the morphism of functors
$D_\QCoh(\mathcal{O}_Y) \to D_\QCoh(\mathcal{O}_{X'})$
given by the composition
\begin{equation}
\label{equation-base-change-map}
L(g')^* \circ a \to
L(g')^* \circ a \circ Rg_* \circ Lg^* \leftarrow
L(g')^* \circ Rg'_* \circ a' \circ Lg^* \to a' \circ Lg^*
\end{equation}
The first arrow comes from the adjunction map $\text{id} \to Rg_* Lg^*$
and the last arrow from the adjunction map $L(g')^*Rg'_* \to \text{id}$.
We need the assumption on Tor independence to invert the arrow
in the middle, see Lemma \href{https://stacks.math.columbia.edu/tag/0A9K}{lemma flat precompose pus (Tag 0A9K)}.
Alternatively, we can think of (\ref{equation-base-change-map}) by
adjointness of $L(g')^*$ and $R(g')_*$ as a natural transformation
$$
a \to a \circ Rg_* \circ Lg^* \leftarrow Rg'_* \circ a' \circ Lg^*
$$
were again the second arrow is invertible. If $M \in D_\QCoh(\mathcal{O}_X)$
and $K \in D_\QCoh(\mathcal{O}_Y)$
then on Yoneda functors this map is given by
\begin{align*}
\Hom_X(M, a(K))
& =
\Hom_Y(Rf_*M, K) \\
& \to
\Hom_Y(Rf_*M, Rg_* Lg^*K) \\
& =
\Hom_{Y'}(Lg^*Rf_*M, Lg^*K) \\
& \leftarrow
\Hom_{Y'}(Rf'_* L(g')^*M, Lg^*K) \\
& =
\Hom_{X'}(L(g')^*M, a'(Lg^*K)) \\
& =
\Hom_X(M, Rg'_*a'(Lg^*K))
\end{align*}
(were the arrow pointing left is invertible by the base
change theorem given in
Derived Categories of Schemes, Lemma \href{https://stacks.math.columbia.edu/tag/08IB}{lemma compare base change (Tag 08IB)})
which makes things a little bit more explicit.

\begin{lemma}
\label{lemma-more-base-change}
In diagram (\ref{equation-base-change}) assume
\begin{enumerate}
\item $g : Y' \to Y$ is a morphism of affine schemes,
\item $f : X \to Y$ is proper, and
\item $f$ and $g$ are Tor independent.
\end{enumerate}
Then the base change map (\ref{equation-base-change-map}) induces an
isomorphism
$$
L(g')^*a(K) \longrightarrow a'(Lg^*K)
$$
in the following cases
\begin{enumerate}
\item for all $K \in D_\QCoh(\mathcal{O}_Y)$ if $f$
is flat of finite presentation,
\item for all $K \in D_\QCoh(\mathcal{O}_Y)$ if $f$
is perfect and $Y$ Noetherian,
\item for $K \in D_\QCoh^+(\mathcal{O}_Y)$ if $g$ has finite Tor dimension
and $Y$ Noetherian.
\end{enumerate}
\end{lemma}

\begin{proof}
Write $Y = \Spec(A)$ and $Y' = \Spec(A')$. As a base change of an affine
morphism, the morphism $g'$ is affine. Let $M$ be a perfect generator
for $D_\QCoh(\mathcal{O}_X)$, see Derived Categories of Schemes, Theorem
\href{https://stacks.math.columbia.edu/tag/09IS}{theorem bondal van den Bergh (Tag 09IS)}. Then $L(g')^*M$ is a
generator for $D_\QCoh(\mathcal{O}_{X'})$, see
Derived Categories of Schemes, Remark \href{https://stacks.math.columbia.edu/tag/0BQT}{remark pullback generator (Tag 0BQT)}.
Hence it suffices to show that (\ref{equation-base-change-map})
induces an isomorphism
\begin{equation}
\label{equation-iso}
R\Hom_{X'}(L(g')^*M, L(g')^*a(K))
\longrightarrow
R\Hom_{X'}(L(g')^*M, a'(Lg^*K))
\end{equation}
of global hom complexes, see
Cohomology, Section \href{https://stacks.math.columbia.edu/tag/0B6A}{section global RHom (Tag 0B6A)},
as this will imply the cone of $L(g')^*a(K) \to a'(Lg^*K)$
is zero.
The structure of the proof is as follows: we will first show that
these Hom complexes are isomorphic and in the last part of the proof
we will show that the isomorphism is induced by (\ref{equation-iso}).

\medskip\noindent
The left hand side. Because $M$ is perfect, the canonical map
$$
R\Hom_X(M, a(K)) \otimes^\mathbf{L}_A A'
\longrightarrow
R\Hom_{X'}(L(g')^*M, L(g')^*a(K))
$$
is an isomorphism by Derived Categories of Schemes, Lemma
\href{https://stacks.math.columbia.edu/tag/0AA7}{lemma affine morphism and hom out of perfect (Tag 0AA7)}.
We can combine this with the isomorphism
$R\Hom_Y(Rf_*M, K) = R\Hom_X(M, a(K))$
of Lemma \href{https://stacks.math.columbia.edu/tag/0B6I}{lemma iso global hom (Tag 0B6I)}
to get that the left hand side equals
$R\Hom_Y(Rf_*M, K) \otimes^\mathbf{L}_A A'$.

\medskip\noindent
The right hand side. Here we first use the isomorphism
$$
R\Hom_{X'}(L(g')^*M, a'(Lg^*K)) = R\Hom_{Y'}(Rf'_*L(g')^*M, Lg^*K)
$$
of Lemma \href{https://stacks.math.columbia.edu/tag/0B6I}{lemma iso global hom (Tag 0B6I)}. Then we use the base change
map $Lg^*Rf_*M \to Rf'_*L(g')^*M$ is an isomorphism by
Derived Categories of Schemes, Lemma \href{https://stacks.math.columbia.edu/tag/08IB}{lemma compare base change (Tag 08IB)}.
Hence we may rewrite this as $R\Hom_{Y'}(Lg^*Rf_*M, Lg^*K)$.
Since $Y$, $Y'$ are affine and $K$, $Rf_*M$ are in $D_\QCoh(\mathcal{O}_Y)$
(Derived Categories of Schemes, Lemma
\href{https://stacks.math.columbia.edu/tag/08D5}{lemma quasi coherence direct image (Tag 08D5)})
we have a canonical map
$$
\beta :
R\Hom_Y(Rf_*M, K) \otimes^\mathbf{L}_A A'
\longrightarrow
R\Hom_{Y'}(Lg^*Rf_*M, Lg^*K)
$$
in $D(A')$. This is the arrow
More on Algebra, Equation (\href{https://stacks.math.columbia.edu/tag/0E1X}{equation base change RHom (Tag 0E1X)})
where we have used Derived Categories of Schemes, Lemmas
\href{https://stacks.math.columbia.edu/tag/06Z0}{lemma affine compare bounded (Tag 06Z0)} and
\href{https://stacks.math.columbia.edu/tag/0A6H}{lemma quasi coherence internal hom (Tag 0A6H)}
to translate back and forth into algebra.
\begin{enumerate}
\item If $f$ is flat and of finite presentation, the complex $Rf_*M$
is perfect on $Y$ by Derived Categories of Schemes, Lemma
\href{https://stacks.math.columbia.edu/tag/0B91}{lemma flat proper perfect direct image general (Tag 0B91)}
and $\beta$ is an isomorphism by
More on Algebra, Lemma \href{https://stacks.math.columbia.edu/tag/0A6A}{lemma base change RHom (Tag 0A6A)} part (1).
\item If $f$ is perfect and $Y$ Noetherian, the complex $Rf_*M$
is perfect on $Y$ by More on Morphisms, Lemma
\href{https://stacks.math.columbia.edu/tag/0B6G}{lemma perfect proper perfect direct image (Tag 0B6G)}
and $\beta$ is an isomorphism as before.
\item If $g$ has finite tor dimension and $Y$ is Noetherian,
the complex $Rf_*M$ is pseudo-coherent on $Y$
(Derived Categories of Schemes, Lemmas
\href{https://stacks.math.columbia.edu/tag/08E2}{lemma direct image coherent (Tag 08E2)} and
\href{https://stacks.math.columbia.edu/tag/08E8}{lemma identify pseudo coherent noetherian (Tag 08E8)})
and $\beta$ is an isomorphism by
More on Algebra, Lemma \href{https://stacks.math.columbia.edu/tag/0A6A}{lemma base change RHom (Tag 0A6A)} part (4).
\end{enumerate}
We conclude that we obtain the same answer as in the previous paragraph.

\medskip\noindent
In the rest of the proof we show that the identifications of
the left and right hand side of (\ref{equation-iso})
given in the second and third paragraph are in fact given by
(\ref{equation-iso}). To make our formulas manageable
we will use $(-, -)_X = R\Hom_X(-, -)$, use $- \otimes A'$
in stead of $- \otimes_A^\mathbf{L} A'$, and we will abbreviate
$g^* = Lg^*$ and $f_* = Rf_*$. Consider the following
commutative diagram
$$
\xymatrix{
((g')^*M, (g')^*a(K))_{X'} \ar[d] &
(M, a(K))_X \otimes A' \ar[l]^-\alpha \ar[d] &
(f_*M, K)_Y \otimes A' \ar@{=}[l] \ar[d] \\
((g')^*M, (g')^*a(g_*g^*K))_{X'} &
(M, a(g_*g^*K))_X \otimes A' \ar[l]^-\alpha &
(f_*M, g_*g^*K)_Y \otimes A' \ar@{=}[l] \ar@/_4pc/[dd]_{\mu'} \\
((g')^*M, (g')^*g'_*a'(g^*K))_{X'} \ar[u] \ar[d] &
(M, g'_*a'(g^*K))_X \otimes A' \ar[u] \ar[l]^-\alpha \ar[ld]^\mu &
(f_*M, K) \otimes A' \ar[d]^\beta \\
((g')^*M, a'(g^*K))_{X'} &
(f'_*(g')^*M, g^*K)_{Y'} \ar@{=}[l] \ar[r] &
(g^*f_*M, g^*K)_{Y'}
}
$$
The arrows labeled $\alpha$ are the maps from
Derived Categories of Schemes, Lemma
\href{https://stacks.math.columbia.edu/tag/0AA7}{lemma affine morphism and hom out of perfect (Tag 0AA7)}
for the diagram with corners $X', X, Y', Y$.
The upper part of the diagram is commutative as the horizontal arrows are
functorial in the entries.
The middle vertical arrows come from the invertible transformation
$g'_* \circ a' \to a \circ g_*$  of Lemma \href{https://stacks.math.columbia.edu/tag/0A9K}{lemma flat precompose pus (Tag 0A9K)}
and therefore the middle square is commutative.
Going down the left hand side is (\ref{equation-iso}).
The upper horizontal arrows provide the identifications used in the
second paragraph of the proof.
The lower horizontal arrows including $\beta$ provide the identifications
used in the third paragraph of the proof. Given $E \in D(A)$,
$E' \in D(A')$, and $c : E \to E'$ in $D(A)$ we will denote
$\mu_c : E \otimes A' \to E'$ the map induced by $c$
and the adjointness of restriction and base change;
if $c$ is clear we write $\mu = \mu_c$, i.e., we
drop $c$ from the notation. The map $\mu$ in the diagram is of this
form with $c$ given by the identification
$(M, g'_*a(g^*K))_X = ((g')^*M, a'(g^*K))_{X'}$
; the triangle involving $\mu$ is commutative by
Derived Categories of Schemes, Remark \href{https://stacks.math.columbia.edu/tag/0BZA}{remark multiplication map (Tag 0BZA)}.

\medskip\noindent
Observe that
$$
\xymatrix{
(M, a(g_*g^*K))_X &
(f_*M, g_* g^*K)_Y \ar@{=}[l] &
(g^*f_*M, g^*K)_{Y'} \ar@{=}[l] \\
(M, g'_* a'(g^*K))_X \ar[u] &
((g')^*M, a'(g^*K))_{X'} \ar@{=}[l] &
(f'_*(g')^*M, g^*K)_{Y'} \ar@{=}[l] \ar[u]
}
$$
is commutative by the very definition of the transformation
$g'_* \circ a' \to a \circ g_*$. Letting $\mu'$ be as above
corresponding to the identification
$(f_*M, g_*g^*K)_X = (g^*f_*M, g^*K)_{Y'}$, then the
hexagon commutes as well. Thus it suffices to show that
$\beta$ is equal to the composition of
$(f_*M, K)_Y \otimes A' \to (f_*M, g_*g^*K)_X \otimes A'$
and $\mu'$. To do this, it suffices to prove the two induced maps
$(f_*M, K)_Y \to (g^*f_*M, g^*K)_{Y'}$ are the same.
In other words, it suffices to show the diagram
$$
\xymatrix{
R\Hom_A(E, K) \ar[rr]_{\text{induced by }\beta} \ar[rd] & &
R\Hom_{A'}(E \otimes_A^\mathbf{L} A', K \otimes_A^\mathbf{L} A') \\
& R\Hom_A(E, K \otimes_A^\mathbf{L} A') \ar[ru]
}
$$
commutes for all $E, K \in D(A)$. Since this is how $\beta$ is constructed in
More on Algebra, Section \href{https://stacks.math.columbia.edu/tag/0E1V}{section base change RHom (Tag 0E1V)}
the proof is complete.
\end{proof}
\end{document}
